% Generated draft: requires mathematical and source review.
\documentclass[11pt]{article}
\usepackage[margin=25mm]{geometry}
\usepackage{fontspec}
\setmainfont{Latin Modern Roman}
\usepackage{amsmath,amssymb,mathtools,unicode-math}
\setmathfont{Latin Modern Math}
\usepackage{enumitem,needspace,xurl,hyperref}
\setlength{\parindent}{0pt}
\setlength{\parskip}{5pt}
\setlength{\emergencystretch}{3em}
\newcommand{\N}{\mathbb{N}}
\newcommand{\Z}{\mathbb{Z}}
\newcommand{\Q}{\mathbb{Q}}
\newcommand{\R}{\mathbb{R}}
\newcommand{\C}{\mathbb{C}}
\newcommand{\F}{\mathbb{F}}
\newcommand{\A}{\mathbb{A}}
\newcommand{\PP}{\mathbb{P}}
\newcommand{\E}{\mathbb{E}}
\newcommand{\eps}{\varepsilon}
\newcommand{\dd}{\,\mathrm{d}}
\begin{document}
\section*{UnsolvedMath ID: 800107}
\textbf{Source problem code:} catalogue entry
\subsection{Definitions and mathematical statement}
Let \(S^1=\{z\in\C:|z|=1\}\). In the usual topological interpretation, an action of a group \(G\) on \(S^1\) is a homomorphism \(\rho:G\to\operatorname{Homeo}(S^1)\), where \(\operatorname{Homeo}(S^1)\) is the group of homeomorphisms of the circle. For a circle lamination \(\Lambda\), invariance under this action means
\[
\rho(g)\Lambda=\Lambda\qquad\text{for every }g\in G,
\]
where the action on laminations is induced by the action on their endpoints. The supplied record does not define either the precise class of circle laminations intended or the additional conditions making a pair of laminations \emph{veering}; these technical meanings must be recovered from the source before a complete formalization is possible.

\textbf{Classification question, with unresolved source conventions.} Classify all \(3\)-manifolds \(M\) for which there exist an action
\[
\rho:\pi_1(M)\longrightarrow\operatorname{Homeo}(S^1)
\]
and two circle laminations \(\Lambda_1,\Lambda_2\) such that
\[
\rho(g)\Lambda_i=\Lambda_i
\qquad\text{for every }g\in\pi_1(M)\text{ and }i\in\{1,2\},
\]
and \((\Lambda_1,\Lambda_2)\) is a veering pair in the intended source sense.

This formulation uses the usual interpretation of a circle action by homeomorphisms; that convention is not explicitly stated in the supplied excerpt. The source also leaves unspecified the intended class of \(3\)-manifolds, restrictions on their boundary, any orientation or faithfulness requirement on the action, and the equivalence relation under which the classification is sought. The notation \(\pi_1(M)\) follows the source; a basepoint and any connectedness convention must be fixed in a complete formulation.

The supplied background reports that known examples of groups admitting such invariant veering pairs are fundamental groups of \(3\)-manifolds carrying either a pseudo-Anosov flow without perfect fits or a veering triangulation. This is retained as attributed background, not as an exhaustive characterization or an independently verified assertion.
\subsection{Short English statement}
The question asks which \(3\)-manifolds have fundamental groups admitting circle actions that preserve a veering pair of laminations. The supplied excerpt does not define the veering condition or specify the full manifold and action hypotheses, so this statement remains conditional on those source conventions. The registry labels the problem open; that label is supplied metadata, not an independently verified status assessment.
\subsection{Sources}
\begin{itemize}
\item Local UnsolvedMath registry, permanent ID \(800107\), source problem code \texttt{catalogue entry}, title \emph{Foliations and Diffeomorphism Groups --- Question 48}. The supplied record attributes the problem to Nariman and Yazdi, with proposed year \(2025\), and labels its status \texttt{open}.
\item Nariman and Yazdi, \emph{Problem List on Foliations and Diffeomorphism Groups} (2025), Question 48. Supplied source URL: \url{https://arxiv.org/abs/2503.18274}. The record reports access on \texttt{2026-10-16} and extraction by \texttt{arxiv\_latex}; these are supplied provenance details, not a claim of independent access or verification.
\item Supplied upstream revision: \texttt{d371ee917ac9647fbb06cd6fc529deec6f3dda82}. The source excerpt supplies the background concerning pseudo-Anosov flows without perfect fits and veering triangulations, but gives no separate references for those examples.
\end{itemize}
\end{document}
